\documentclass[11pt]{amsart}
\usepackage[T1]{fontenc}
\usepackage{lmodern}
\usepackage{amsmath,amssymb,amsthm,mathtools}
\usepackage{xcolor}
\usepackage[colorlinks=true,linkcolor=blue!55!black,citecolor=blue!55!black,urlcolor=blue!55!black]{hyperref}
\raggedbottom

\newtheorem{theorem}{Theorem}[section]
\newtheorem{lemma}[theorem]{Lemma}
\newtheorem{proposition}[theorem]{Proposition}
\newtheorem{corollary}[theorem]{Corollary}
\theoremstyle{definition}
\newtheorem{definition}[theorem]{Definition}
\newtheorem{remark}[theorem]{Remark}

\title{The Exact Finite Sumset Value \(s(12,38)=65\)}
\author{Zijian Zeng}
\address{Institute of Computer Science and Digital Innovation, UCSI University, Kuala Lumpur, 56000, MALAYSIA}
\email{zijianzeng@foxmail.com; 1002266693@ucsiuniversity.edu.my}
\date{2026-09-01}
\hypersetup{pdftitle={The Exact Finite Sumset Value s(12,38)=65},pdfauthor={Zijian Zeng}}

\begin{document}

\begin{abstract}
For \(s(k,n)=\max\{|A+A|:A\subseteq[1,n],\ |A|=k\}\), we prove the
exact finite value \(s(12,38)=65\).  The lower bound is witnessed by an
explicit integer set.  The upper bound is a finite computational proposition:
a definition-level CNF encoding of the existence of a set with at least 66
distinct sums has a DRUP/DRAT refutation that is replayed by an independent
proof checker.  Exhaustive small-instance comparisons test the encoding.
This result concerns one finite cell and has no claimed asymptotic consequence.
\end{abstract}

\maketitle
\enlargethispage{4pt}

\section{Introduction}

Pikhurko introduced the finite extremal quantity
\[
s(k,n)=\max\{|A+A|:A\subseteq[1,n],\ |A|=k\}
\]
in connection with edge-magic graphs and quasi-Sidon sets
\cite{Pikhurko2006}.  An online computational report dated August 12, 2026,
records exact values through \(n=37\), including \(s(12,37)=64\)
\cite{EPAD840}.  Its linked certificate bundle was not accessible during the
present audit, so that row is background rather than an input to our proof.

We determine only the next specified cell.  Exact-phrase, repository, and
forward-source searches through September 1, 2026, located no earlier public
determination of \(s(12,38)\).  This bounded search statement is not a claim
that an uncatalogued computation cannot exist.

\section{Definitions and reductions}

Write \([1,n]=\{1,\ldots,n\}\) and
\(A+A=\{a+b:a,b\in A\}\); diagonal sums are included.  There are
\(\binom{13}{2}=78\) unordered pairs with repetition for a 12-set, and all
sums lie in \([2,76]\).

Translation preserves the objective.  Every nonempty subset of \([1,38]\)
can therefore be translated to have minimum 1 while remaining in
\([1,38]\).  The upper-bound encoding imposes this normalization.

\section{Main result}

\begin{theorem}\label{thm:main}
Let
\[
s(12,38)=\max_{\substack{A\subseteq\{1,\ldots,38\}\\|A|=12}}|A+A|.
\]
Then \(s(12,38)=65\).
\end{theorem}

\begin{proof}
Let
\[
A=\{1,2,4,8,13,18,21,26,31,35,37,38\}.
\]
Direct integer evaluation gives the 65 sums
\begin{align*}
&2,3,4,5,6,8,9,10,12,14,15,16,17,19,20,21,22,23,25,26,27,28,\\
&29,30,31,32,33,34,35,36,37,38,39,40,41,42,43,44,45,46,47,48,\\
&49,50,51,52,53,55,56,57,58,59,61,62,63,64,66,68,69,70,72,73,\\
&74,75,76.
\end{align*}
Thus \(s(12,38)\ge65\).

For the reverse inequality, Proposition~\ref{prop:certificate} certifies that
no normalized 12-subset has 66 distinct sums.  Translation shows that every
12-subset has a normalized representative, so \(s(12,38)\le65\).
\end{proof}

\section{Exact finite certificate}

\begin{proposition}\label{prop:certificate}
There is no 12-element \(A\subseteq[1,38]\) satisfying \(|A+A|\ge66\).
\end{proposition}

The CNF has selection variables \(x_i\) for \(1\le i\le38\), pair variables
\(z_{ij}\) for \(1\le i<j\le38\), and sum variables \(y_t\) for
\(2\le t\le76\).  Sequential-counter clauses impose
\(\sum_i x_i=12\), and the unit clause \(x_1\) imposes the translation
normalization.  For \(i<j\), three clauses express
\(z_{ij}\leftrightarrow(x_i\wedge x_j)\); for diagonal pairs we use
\(z_{ii}=x_i\).  For each \(t\), clauses express
\[
y_t\longleftrightarrow\bigvee_{\substack{1\le i\le j\le38\\i+j=t}}z_{ij}.
\]
A final sequential counter imposes \(\sum_{t=2}^{76}y_t\ge66\).
Consequently the CNF is satisfiable exactly when the forbidden set exists.

The generated DIMACS instance has 2034 variables and 5419 clauses.  A
proof-producing SAT solver emitted an ASCII DRUP/DRAT refutation.  The
separate \texttt{drat-trim} checker \cite{WetzlerHeuleHunt2014} reported
\texttt{VERIFIED}; its backward check retained 3854 input clauses and 59,494
lemmas and performed 6,690,497 resolution steps.  This proves the proposition.

As a separate attack on encoding errors, transparent enumeration of every
subset for all 55 cells \(1\le k\le n\le10\) was compared with the SAT encoding
at the optimum and the two adjacent thresholds.  All comparisons passed.  The
target CNF was regenerated after the test and was byte-for-byte identical.

\section{Limitations}

The theorem establishes an attained maximum but does not classify all
maximizing 12-sets.  It does not compute another entry of the \(n=38\) row and
does not strengthen any asymptotic quasi-Sidon bound.

\section*{Data, code, and computational disclosure}

The accompanying archive contains the literal witness, an independent integer
evaluator, the definition-to-CNF generator, the DIMACS instance, the DRUP/DRAT
proof, the checker log, and the small-instance cross-check.  The full commands
and software versions are recorded in \texttt{REPRODUCE.md}.  Computational
search was used to discover the witness and refutation; discovery and
certificate verification are separated, and the theorem relies only on the
displayed exact witness and the replayed finite proof.

\bibliographystyle{amsplain}
\bibliography{references}

\end{document}
